%==============================================================================
% ABSTRACT
%==============================================================================

\chapter*{Abstract}
\addcontentsline{toc}{chapter}{Abstract}
\markboth{Abstract}{Abstract}

\noindent
This project presents \nabta, a comprehensive AI-powered smart agriculture platform that integrates computer vision, machine learning, explainable artificial intelligence, and large language models to support intelligent agricultural decision-making. The platform provides an end-to-end ecosystem for automated plant disease diagnosis, crop recommendation, irrigation recommendation, and AI-assisted agricultural consultation through a unified web-based application. By combining multiple artificial intelligence technologies within a single production-ready platform, \nabta aims to improve agricultural productivity, optimize resource utilization, and promote sustainable farming practices.

The plant disease analysis module employs a multi-stage computer vision pipeline. \gls{yolov8} performs real-time object detection to localize infected leaf regions, followed by a \unet segmentation model with an EfficientNet-B0 encoder that isolates leaf tissue from complex backgrounds using pixel-level segmentation. The refined leaf images are then classified by a custom \gls{cnn} trained on a unified dataset containing more than \textbf{120,000} images across \textbf{102} healthy and diseased plant classes. To improve model transparency and user trust, \gls{gradcam} is incorporated to generate visual explanations highlighting the image regions responsible for each prediction.

Beyond disease diagnosis, \nabta integrates intelligent decision-support services for precision agriculture. A Random Forest-based crop recommendation model analyzes soil properties, including Nitrogen (N), Phosphorus (P), Potassium (K), soil type, weather conditions, and environmental factors to recommend the most suitable crops for cultivation. In addition, the platform includes two irrigation intelligence modules: a Random Forest classification model for irrigation decision-making and a Random Forest regression model for estimating crop water requirements, enabling efficient irrigation planning and optimized water resource management.

To further enhance user interaction, the platform incorporates a domain-specific conversational assistant powered by the Cohere \gls{api} and Retrieval-Augmented Generation (\gls{rag}). The assistant provides evidence-based disease treatment recommendations, crop management guidance, irrigation advice, and general agricultural consultations in both Arabic and English. All intelligent services are deployed through a high-performance asynchronous \fastapi backend and presented via a responsive React.js dashboard featuring image analysis, recommendation modules, historical records, analytics, authentication, and user management.

Experimental evaluation demonstrates the effectiveness of the proposed AI modules. The plant disease classification model achieved an overall accuracy of \textbf{95.0\%} across \textbf{102} healthy and diseased plant classes trained on more than \textbf{120,000} images. The YOLOv8x detection model achieved a validation Precision of \textbf{72.9\%}, Recall of \textbf{72.6\%}, and an \gls{map}@0.50 of \textbf{78.8\%}, while the proposed U-Net segmentation model achieved a best mean \gls{iou} of \textbf{94.73\%}. The crop recommendation model achieved a test accuracy of \textbf{85.13\%} with a cross-validation accuracy of \textbf{84.8\%}, while the irrigation decision model achieved a test accuracy of \textbf{98.53\%}. Furthermore, the irrigation water requirement regression model achieved a test coefficient of determination (\(R^2\)) of \textbf{95.22\%} with a mean absolute error of only \textbf{13.47 mm}. Collectively, these results demonstrate the capability of \nabta to deliver accurate plant disease diagnosis, intelligent agricultural recommendations, efficient irrigation management, and scalable AI-powered decision support for modern precision agriculture.

\vspace{1em}
\noindent\textbf{Keywords:} Smart agriculture, plant disease detection, crop recommendation, irrigation recommendation, deep learning, machine learning, Random Forest, YOLOv8, U-Net, convolutional neural networks, Grad-CAM, explainable AI, retrieval-augmented generation, FastAPI, precision agriculture.
